%------------------------------------------------------------------------------%
\documentclass[fleqn,utf8,aspectratio=169,12pt]{beamer}

\usepackage{gaal}

\title{Determinantes}

%------------------------------------------------------------------------------%
\begin{document}

\addtitle

%------------------------------------------------------------------------------%
\section{Determinantes}

%------------------------------------------------------------------------------%
\begin{frame}{Determinantes}
  O determinante é uma função que associa um número real a uma matriz quadrada

  \pause
  \medskip

  Para uma matriz quadrada $A$, o determinante é denotado por
  \[
    \det(A)
  \]
  \pause
  outra notação para o determinante é colocar os elementos entre linhas verticais
  \[
    \det(A) = \abs{A}
  \]
\end{frame}

%------------------------------------------------------------------------------%
\begin{frame}{Determinantes}
  O determinante possui aplicações importantes
  \begin{itemize}[<+->]
%    \setlength{\itemsep}{1em}
    \item verificar se uma matriz é invertível
    \item analisar sistemas lineares
    \item determinar dependência linear
    \item calcular áreas e volumes
    \item estudar transformações lineares
  \end{itemize}
\end{frame}

%------------------------------------------------------------------------------%
\section{Definição}

%------------------------------------------------------------------------------%
\begin{frame}{Determinante de uma matriz $1\times 1$}
  Para uma matriz de ordem 1
  \[
    A =
    \begin{pmatrix}
      a_{11}
    \end{pmatrix}
  \]
  \pause
  definimos
  \[
    \det(A) = a_{11}
  \]
\end{frame}

%------------------------------------------------------------------------------%
\begin{frame}{Determinante de uma matriz $2\times 2$}
  O determinante da matriz $2\times 2$
  \[
    A =
    \begin{pmatrix}
      {\color<7>{magenta} a} & {\color<7>{green} b}   \\
      {\color<7>{green} c}   & {\color<7>{magenta} d}
    \end{pmatrix}
  \]

  \pause
  é o produto dos elementos da \emph{diagonal}
  \pause
  \emph{menos}
  \\ \pause
  o produto dos elementos da \emph{diagonal secundária} (invertida)
  \pause
  \[
    \det(A)
    \pause
    \;=\;
    \begin{vmatrix}
      {\color<7>{magenta} a} & {\color<7>{green} b}   \\
      {\color<7>{green} c}   & {\color<7>{magenta} d}
    \end{vmatrix}
    \pause
    \;=\;
    {\color<7>{magenta}ad} - {\color<7>{green}bc}
  \]

\end{frame}

%------------------------------------------------------------------------------%
%------------------------------------------------------------------------------%
\begin{question}
  Calcule o determinante da matriz
  \[
    A =
    \begin{pmatrix}
      3 & 2 \\
      1 & 4
    \end{pmatrix}
  \]
\end{question}

%------------------------------------------------------------------------------%
\begin{resolution}{Solução}
  \begin{align*}
    \onslide<+->{\det(A)}
    \onslide<+->{& =
      \begin{vmatrix}
        3 & 2 \\
        1 & 4
      \end{vmatrix} \\}
    \onslide<+->{& = a_{11}a_{22} - a_{12}a_{21} \\}
    \onslide<+->{& = 3 \times 4 - 2 \times 1 \\}
    \onslide<+->{& = 12 - 2 \\}
    \onslide<+->{& = 10}
  \end{align*}
\end{resolution}

%------------------------------------------------------------------------------%
\endinput


%------------------------------------------------------------------------------%
\begin{frame}{Determinante de uma matrizes $n\times n$}
  \onslide<+->{Para definir determinantes de matrizes de ordem maior}

  \onslide<+->{Precisamos definir operações conceitos intermediários}
  \begin{itemize}[<+->]
  \item Menor complementar
  \item Cofator
  \end{itemize}
  \onslide<+->{para apresentarmos o \emph{Teorema de Laplace}}
\end{frame}

%------------------------------------------------------------------------------%
\begin{frame}{Menor complementar}
  Dada uma matriz $A=(a_{ij})$ de ordem $n$

  \pause
  O \emph{menor complementar} $M_{ij}$
  \pause
  é o determinante da matriz obtida

  \pause
  eliminando a linha $i$ e a coluna $j$ de $A$
\end{frame}

%------------------------------------------------------------------------------%
%------------------------------------------------------------------------------%
\begin{question}
  Determine o menor complementar associado ao elemento $a_{12}$
  da matriz
  \[
    A =
    \begin{bmatrix}
      1 & 2 & 3 \\
      4 & 5 & 6 \\
      7 & 8 & 9
    \end{bmatrix}
  \]
\end{question}

%------------------------------------------------------------------------------%
\begin{resolution}{Solução}
  \onslide<+->{Obtemos o menor complementar eliminando a primeira linha e a segunda coluna}
  \[
    \onslide<1->{A =}
    \begin{bNiceMatrix}
      \onslide<1->{1 & 2 & 3} \\
      \onslide<1->{4 & 5 & 6} \\
      \onslide<1->{7 & 8 & 9}
    \CodeAfter
      \onslide<+->{\tikz \draw[blue, thick] (1.5-|1) -- (1.5-|last); }
      \onslide<+->{\tikz \draw[blue, thick] (1-|2.5) -- (last-|2.5); }
    \end{bNiceMatrix}
  \]
  \[
    \onslide<+->{M_{12}}
    \onslide<+->{\;=\;
    \det
    \begin{bmatrix}
      4 & 6 \\
      7 & 9
    \end{bmatrix}}
    \onslide<+->{\;=\;
    \begin{vmatrix}
      4 & 6 \\
      7 & 9
    \end{vmatrix}}
    \onslide<+->{\;=\; 4 \times 9 - 6 \times 7}
    \onslide<+->{\;=\; -7}
  \]
\end{resolution}

%------------------------------------------------------------------------------%
\endinput






%------------------------------------------------------------------------------%
\begin{frame}{Cofator}
  O \emph{cofator} associado ao elemento $a_{ij}$ é definido por
  \pause
  \[
    C_{ij} = (-1)^{i+j }M_{ij}
  \]

  \pause
  \bigskip
  Isto é
  \[
    \pause
    C_{ij} = \hpm M_{ij}
    \qquad
    \text{ quando } \;i+j\; \text{ é par}
  \]
  \[
    \pause
    C_{ij} = -M_{ij}
    \qquad
    \text{ quando } \;i+j\; \text{ é ímpar}
  \]
\end{frame}

%------------------------------------------------------------------------------%
%------------------------------------------------------------------------------%
\begin{question}
  Determine o cofator associado ao elemento $a_{12}$
  da matriz
  \[
    A =
    \begin{bmatrix}
      1 & 2 & 3 \\
      4 & 5 & 6 \\
      7 & 8 & 9
    \end{bmatrix}
  \]
\end{question}

%------------------------------------------------------------------------------%
\begin{resolution}{Solução}
  \[
    \onslide<+->{C_{12}}
    \onslide<+->{\;=\; (-1)^{1+2}M_{12}}
    \onslide<+->{\;=\; -M_{12}}
    \onslide<+->{\;=\;
    -\det
    \begin{bmatrix}
      4 & 6 \\
      7 & 9
    \end{bmatrix}}
    \onslide<+->{\;=\;
      -\begin{vmatrix}
        4 & 6 \\
        7 & 9
      \end{vmatrix}}
  \]
  \[
    \hphantom{C_{12}}
    \onslide<+->{\;=\; -(4 \times 9 - 6 \times 7)}
    \onslide<+->{\;=\; 7}
  \]
\end{resolution}

%------------------------------------------------------------------------------%
\endinput






%------------------------------------------------------------------------------%
\section{Teorema de Laplace}

%------------------------------------------------------------------------------%
\begin{frame}{Teorema de Laplace}
  \emph{Teorema de Laplace}

  \pause
  permite calcular o determinante de uma matriz quadrada

  \pause
  expandindo-o por qualquer linha ou qualquer coluna

  \pause
  \bigskip
  \bigskip
  A escolha da linha ou coluna é livre
\end{frame}

%------------------------------------------------------------------------------%
\begin{frame}{Teorema de Laplace -- Expansão por Linha}
  Expandindo pela linha $i$
  \[
    \det(A) \;=\;
    \sum_{\emph{j}=1}^{n} a_{i\emph{j}} C_{i\emph{j}}
  \]
  \[
    C_{ij} = (-1)^{i+j }M_{ij}
  \]
\end{frame}

%------------------------------------------------------------------------------%
\begin{frame}{Teorema de Laplace -- Expansão por Coluna}
  Expandindo pela coluna $j$
  \[
    \det(A) \;=\;
    \sum_{\vphantom{j}\emph{i}=1}^{n} a_{\emph{i}j} C_{\emph{i}j}
  \]
  \[
    C_{ij} = (-1)^{i+j }M_{ij}
  \]
\end{frame}

%------------------------------------------------------------------------------%
%------------------------------------------------------------------------------%
\begin{question}
  Calcule o determinante da matriz usando Teorema de Laplace
  \[
    A =
    \begin{pmatrix}
      1 & 2 & 3 \\
      0 & 4 & 5 \\
      0 & 6 & 7
    \end{pmatrix}
  \]
\end{question}

%------------------------------------------------------------------------------%
\begin{resolution}{Solução}
  \onslide<+->{\[
    A =
    \begin{pmatrix}
      1 & 2 & 3 \\
      0 & 4 & 5 \\
      0 & 6 & 7
    \end{pmatrix}
  \]}%
  \onslide<+->{Podemos expandir o determinante por linha ou coluna}

  \onslide<+->{Escolhendo a primeira \emph{coluna} temos mais elementos nulos}
  \begin{align*}
    \onslide<+->{\det(A)}
    \onslide<+->{\;=\; \sum_{\emph{i}=1}^{3} a_{\emph{i}1} C_{\emph{i}1}}
    \onslide<+->{&\;=\; a_{11} C_{11}}
    \onslide<+->{+ a_{21} C_{21}}
    \onslide<+->{+ a_{31} C_{31}} \\
    \onslide<+->{&\;=\; 1 \; C_{11} + 0 \; C_{21} + 0 \; C_{31}}
  \end{align*}
\end{resolution}

%------------------------------------------------------------------------------%
\begin{resolution}{Solução}
  \begin{align*}
    \onslide<+->{\det(A)}
    \onslide<+->{& = a_{11} C_{11} + a_{21} C_{21} + a_{31} C_{31}     \\}
    \onslide<+->{& = 1 \; C_{11} + 0 \; C_{21} + 0 \; C_{31}           \\}
    \onslide<+->{& = C_{11}                                            \\}
    \onslide<+->{& = (-1)^{1+1} M_{11}                                 \\}
    \onslide<+->{& = \begin{vmatrix}
          4 & 5 \\
          6 & 7
        \end{vmatrix} \\}
    \onslide<+->{& = 4 \times 7 - 5 \times 6}
    \onslide<+->{\;=\; 28 - 30}
    \onslide<+->{\;=\; -2}
  \end{align*}
\end{resolution}

%------------------------------------------------------------------------------%
\endinput

\[
\det(A)
\]
\end{frame}

\begin{frame}{Exemplo -- Laplace por uma linha}
Continuando,
\[
\det(A)=.
\]
Portanto,
\[
\det(A)=
\]
\medskip

Assim,
\[
\boxed{\det(A)=-2}
\]
\medskip

\textbf{Estratégia:} procure uma linha ou coluna que contenha
muitos zeros para simplificar a expansão de Laplace.
\end{frame}

%------------------------------------------------------------------------------%
\begin{question}
  Calcule o determinante da matriz usando Teorema de Laplace
  \[
    A =
    \begin{pmatrix}
      2 & 0 & 1 \\
      3 & 4 & 0 \\
      1 & 2 & 5
    \end{pmatrix}
  \]
\end{question}

%------------------------------------------------------------------------------%
\begin{resolution}{Solução}
  \onslide<+->{\[
    A =
    \begin{pmatrix}
      2 & 0 & 1 \\
      3 & 4 & 0 \\
      1 & 2 & 5
    \end{pmatrix}
  \]}
  \onslide<+->{A segunda linha contém um zero}
  \begin{align*}
    \onslide<+->{\det(A)}
    \onslide<+->{\;=\; \sum_{\emph{j}=1}^{3} a_{2\emph{j}} C_{2\emph{j}}}
    \onslide<+->{&\;=\; a_{21} C_{21}}
    \onslide<+->{+ a_{22} C_{22}}
    \onslide<+->{+ a_{23} C_{23}} \\
    \onslide<+->{&\;=\; 3 \; C_{21} + 4 \; C_{22} + 0 \; C_{23}} \\
    \onslide<+->{&\;=\; 3 \; C_{21} + 4 \; C_{22}}
  \end{align*}
\end{resolution}

%------------------------------------------------------------------------------%
\begin{resolution}{Calculando os Cofatores}
\begin{columns}[t]
\column{0.35\linewidth}
  \begin{align*}
    \onslide<+->{C_{21}}
    \onslide<+->{& = (-1)^{2+1} M_{21} \\}
    \onslide<+->{& = -M_{21} \\}
    \onslide<+->{& = - \begin{vmatrix}
          0 & 1 \\
          2 & 5
        \end{vmatrix} \\}
    \onslide<+->{& = -(0 \times 5 - 1 \times 2) \\}
    \onslide<+->{& = 2}
  \end{align*}
\column{0.5\linewidth}
  \begin{align*}
    \onslide<+->{C_{22}}
    \onslide<+->{& = (-1)^{2+2} M_{22} \\}
    \onslide<+->{& = M_{22} \\}
    \onslide<+->{& = \begin{vmatrix}
          2 & 1 \\
          1 & 5
        \end{vmatrix} \\}
    \onslide<+->{& = 2 \times 5 - 1 \times 1 \\}
    \onslide<+->{& = 9}
  \end{align*}
\end{columns}
\end{resolution}

%------------------------------------------------------------------------------%
\begin{resolution}{Calculando os Cofatores}
  \begin{align*}
    \onslide<+->{\det(A)}
    \onslide<+->{& = 3 \; C_{21} + 4 \; C_{22}} \\
    \onslide<+->{& = 3 \times 2 + 4 \times 9}   \\
    \onslide<+->{& = 6 + 36}                    \\
    \onslide<+->{& = 42}
  \end{align*}
\end{resolution}

%------------------------------------------------------------------------------%
\endinput


%------------------------------------------------------------------------------%
\section{Regra de Sarrus}

%------------------------------------------------------------------------------%
\begin{frame}{Regra de Sarrus}
  \emph{APENAS} para matrizes $3\times 3$

  \pause
  existe um procedimento conhecido como \emph{Regra de Sarrus}

  \pause
  Para
  \[
    A=
    \begin{pmatrix}
      a & b & c \\
      d & e & f \\
      g & h & i
    \end{pmatrix}
  \]
  \pause
  \[
    \det(A) = aei + bfg + cdh - ceg - bdi - afh
  \]
\end{frame}

%------------------------------------------------------------------------------%
\begin{frame}{Diagrama da Regra de Sarrus}
  \[
    \begin{bNiceMatrix}
      a & b & c & a & b \\
      d & e & f & d & e \\
      g & h & i & g & h
    \CodeAfter
      \tikz \draw[black,   thick] (1-|4) -- (4-|4);
      \onslide<2->{\tikz \draw[blue,    thick] (1-|1) -- (4-|4);}
      \onslide<3->{\tikz \draw[blue,    thick] (1-|2) -- (4-|5);}
      \onslide<4->{\tikz \draw[blue,    thick] (1-|3) -- (4-|6);}
      \onslide<5->{\tikz \draw[magenta, thick] (1-|6) -- (4-|3);}
      \onslide<6->{\tikz \draw[magenta, thick] (1-|5) -- (4-|2);}
      \onslide<7->{\tikz \draw[magenta, thick] (1-|4) -- (4-|1);}
    \end{bNiceMatrix}
  \]

  \bigskip
  \[
    \onslide<1->{\det(A)}
    \onslide<2->{\;=\; aei}
    \onslide<3->{\;+\; bfg}
    \onslide<4->{\;+\; cdh}
    \onslide<5->{\;-\; ceg}
    \onslide<6->{\;-\; bdi}
    \onslide<7->{\;-\; afh}
  \]
\end{frame}

%------------------------------------------------------------------------------%
%------------------------------------------------------------------------------%
\begin{question}
Calcule o determinante da matriz pela Regra de Sarrus
  \[
    A =
    \begin{pmatrix}
      1 & 2 & 3 \\
      2 & 1 & 4 \\
      3 & 0 & 1
    \end{pmatrix}
  \]
\end{question}

%------------------------------------------------------------------------------%
\begin{resolution}{Solução}
  \[
    \begin{bNiceMatrix}
      1 & 2 & 3 & 1 & 2 \\
      2 & 1 & 4 & 2 & 1 \\
      3 & 0 & 1 & 3 & 0
    \CodeAfter
      \tikz \draw[black,   thick] (1-|4) -- (4-|4);
      \onslide<2->{\tikz \draw[blue,    thick] (1-|1) -- (4-|4);}
      \onslide<3->{\tikz \draw[blue,    thick] (1-|2) -- (4-|5);}
      \onslide<4->{\tikz \draw[blue,    thick] (1-|3) -- (4-|6);}
      \onslide<5->{\tikz \draw[magenta, thick] (1-|6) -- (4-|3);}
      \onslide<6->{\tikz \draw[magenta, thick] (1-|5) -- (4-|2);}
      \onslide<7->{\tikz \draw[magenta, thick] (1-|4) -- (4-|1);}
    \end{bNiceMatrix}
  \]
  \begin{align*}
    \onslide<1->{\det(A)}
    \onslide<2->{& = 1 \!\times\! 1 \!\times\! 1}
    \onslide<3->{\;+ 2 \!\times\! 4 \!\times\! 3}
    \onslide<4->{\;+ 3 \!\times\! 2 \!\times\! 0}
    \onslide<5->{\;- 3 \!\times\! 1 \!\times\! 3}
    \onslide<6->{\;- 2 \!\times\! 2 \!\times\! 1}
    \onslide<7->{\;- 1 \!\times\! 4 \!\times\! 0} \\
    \onslide<8->{& = 1 + 24 + 0 - 9 - 4 - 0} \\
    \onslide<9->{& = 25 - 13} \\
    \onslide<10->{& = 12}
  \end{align*}
\end{resolution}

%------------------------------------------------------------------------------%
\endinput





%------------------------------------------------------------------------------%
%\section{Lista Mínima}
%
%%------------------------------------------------------------------------------%
%\begin{frame}{Lista Mínima}
%  \begin{enumerate}
%    \item
%  \end{enumerate}
%
%  \vfill
%  Atenção: A prova é baseada no livro, não nas apresentações
%\end{frame}

%------------------------------------------------------------------------------%
\end{document}
%------------------------------------------------------------------------------%
